% =====================================================================
% §2.3 — La posizione di lettura (esempio: ex2_pointer)
%
% Pattern: domanda musicale -> diff YAML (un solo blocco) ->
%          fallimento della waveform -> PRIMA figura-partitura con
%          legenda minima -> meccanismo -> artefatto del freeze ->
%          ponte a §2.4
%
% TODO prima della submission:
%   [ ] linerange del listato: allineare alle righe reali di
%       examples/ex2_pointer/ex2_pointer.yml (mostrare solo il blocco pointer)
%   [ ] verificare nome figura: ex2_pointer_score.pdf
%   [ ] verificare il punto di congelamento (~1,3 s) dopo la taratura
%       definitiva dell'envelope, e che cada su vocale
%   [ ] \cite{Lippe1993cim}: allineare alla chiave reale in refs.bib
% =====================================================================

% ----------------------------------------------------------------
\subsection{La posizione di lettura}\label{sec:pointer}

Nella granulazione di campioni i tempi in gioco sono due: il tempo del
suono registrato e il tempo con cui lo si percorre. Nello stream minimo
coincidono --- è parte della copia fedele --- e lo scostamento
successivo li separa: che cosa accade quando il tempo di lettura si
stacca dal tempo del materiale?

    \lstinputlisting[
      linerange={40-42}, % TODO: allineare al file reale
      language=yaml,
      basicstyle=\ttfamily\fontsize{7.5}{9}\selectfont,
      label=lst:pointer,
      caption={L'unica chiave aggiunta allo stream minimo
      (\texttt{ex2\_pointer.yml}).}
    ]{examples/ex2_pointer/ex2_pointer.yml}

Il diff è un solo blocco (Listato~\ref{lst:pointer}): un envelope sulla
velocità di lettura. Per i primi quattro decimi di secondo
\texttt{speed\_ratio} vale 1 e lo stream \emph{è} la copia fedele di
\S\ref{sec:stream-minimo}; poi la velocità decresce fino a zero e la
testina si arresta a circa 1,3~s nel campione, su una vocale. All'ascolto
è il gesto granulare per antonomasia: la parola che rallenta e si fa
vocale tenuta, il \emph{time-stretch} che degenera in congelamento.

Ma di questa trasformazione la forma d'onda --- lo strumento di verifica
delle Figg.~\ref{fig:stream-minimo-waveform}
e~\ref{fig:griglia-waveform} --- non dà conto: l'energia resta piena e
continua, perché ciò che è cambiato non è \emph{quanta} energia scorre ma
\emph{quale punto} del materiale la genera. Serve una rappresentazione
che mostri la lettura, ed è la partitura grafica che il sistema genera a
ogni rendering (Fig.~\ref{fig:pointer-score}): in ascissa il tempo del
brano, in ordinata la posizione di lettura nel buffer sorgente, ogni
segmento un grano. La traiettoria dichiarata nello YAML vi appare
letteralmente disegnata: la diagonale della lettura naturale che si piega
fino all'orizzontale del congelamento. Alla notazione, qui ridotta alla
sua legenda minima, è dedicata la \S\ref{sec:partitura}.

\begin{figure}[h]
    \includegraphics[width=\linewidth]{examples/ex2_pointer/ex2_pointer_score}\\
    \captionof{figure}{Partitura grafica del rendering di
    \texttt{ex2\_pointer}: ascisse = tempo del brano, ordinate =
    posizione di lettura nel buffer sorgente (0 = inizio, 1 = fine del
    campione); ogni segmento un grano. La diagonale (lettura a velocità
    naturale) decelera fino all'orizzontale (congelamento).}
    \label{fig:pointer-score}
\end{figure}

Il meccanismo è elementare: la posizione di lettura è l'integrale della
velocità, e \texttt{speed\_ratio} ammette qualunque valore --- 1 è la
lettura naturale, 0 il congelamento, i valori negativi la lettura
retrograda. È su questo asse che la granulazione di campioni colloca il
proprio parametro espressivo dominante~\cite{Lippe1993cim}, e l'esempio
mostra perché: stesso materiale, stessa griglia di emissione, e un tempo
musicale interamente altro. Si noti che la specifica è ancora del tutto
deterministica: due rendering coincidono al campione.

Proprio questo determinismo, però, ha un costo udibile. Nel tratto
congelato i grani sono identici e nascono su una griglia metronomica: il
risultato è periodico, un ronzio meccanico che tradisce il granulatore.
Come si scioglie questa rigidità senza rinunciare alla traiettoria
disegnata?



